\section{Introduction}
\label{sec:intro}

Multimodal learning has matured on settings where modalities are structurally synchronized: video frames align with audio waveforms, captions are written with the image present, and transcripts are anchored to utterance boundaries~\cite{baltrusaitis2019multimodal,tsai2019multimodal}.
Real systems are often less forgiving.
External context can arrive late, sparsely, and with variable usefulness relative to a continuously evolving stream.
We refer to this challenge as \emph{asynchronous alignment}: the problem of fusing a dense primary signal with sporadic secondary context whose predictive value depends on when it arrives.
Despite its practical importance, asynchronous alignment remains underexplored because many standard benchmarks hide it through structural synchrony and because temporal validity is usually treated as metadata rather than as a modeling target~\cite{baltrusaitis2019multimodal,dasgupta2018hyte,liu2021temporal}.
In this paper we focus on one especially important subcase, \emph{event-conditioned asynchronous fusion}, where a dense primary stream must be updated using delayed textual context emitted by external events.
High-frequency cryptocurrency markets provide a particularly strict stress test for this problem: price time-series evolve under a dense forward-only clock, while web intelligence---news, alerts, and sentiment cascades---arrives intermittently with delays that can materially change its value.
A news item published 8 minutes after a liquidity shock carries very different signal than one published 3 hours later, after the market has largely adjusted.
This \emph{modality lag} ($\tau_{\mathrm{lag}}$) is therefore not noise to normalize away, but a first-class property of asynchronous multimodal prediction.

Existing multimodal architectures handle asynchronous signals inadequately~\cite{baltrusaitis2019multimodal,tsai2019multimodal}.
Methods requiring strict synchronization~\cite{sawhney2020multimodal,qin2019mdrm} discard most available intelligence, as modalities rarely co-occur within narrow windows.
Window-aggregation approaches~\cite{xu2018stocknet,ding2015deepevent} treat fresh and stale signals as interchangeable, conflating pre- versus post-event information.
Large language models~\cite{lopez2023fingpt,wu2023bloomberggpt} excel at language understanding but lack architectures for asynchronous alignment.
The result is a systematic blind spot: modality lag is rarely modeled as an architectural input.

We therefore frame asynchronous alignment as a core financial multimodal learning problem.
We introduce \dataset{}, an evaluation corpus that pairs high-frequency price sequences with lagged web intelligence across Bitcoin, Ethereum, and Solana.
\dataset{} records $\tau_{\mathrm{lag}}$ explicitly for lag-conditioned modeling and stratified analysis.
Its primary real-news split contains 27{,}914 pooled BTC/ETH/SOL samples from CryptoCompare over December 2025--March 2026, with a controlled synthetic split used only for development checks.
Cryptocurrency markets serve as the target financial domain in this paper and as a strict stress test for asynchronous fusion.

We propose \CGCMA{}, a conditional cross-modal gating architecture for asynchronous prediction.
Text first attends over the full price sequence to produce grounded context, after which a per-dimension residual gate decides how much of that context to inject based on price--text agreement, web features, and $\tau_{\mathrm{lag}}$ (see Section~\ref{sec:method}).
The gate can close automatically when external context is stale or contradictory, degrading gracefully toward unimodal prediction.
We evaluate \CGCMA{} on \dataset{} against price-only, text-only, representative generic multimodal baselines~\cite{xu2018stocknet,tsai2019multimodal,zadeh2017tensor}, and task-specific lag-aware controls under a shared protocol.
The paper's central claim is therefore not that cryptocurrency forecasting is solved, but that event-conditioned asynchronous fusion can be operationalized as a concrete multimodal setting with measurable freshness structure, controlled evaluation, and a plausible architecture for lag-aware trust.

Our contributions are three-fold:
\begin{enumerate}
  \item \textbf{We define and instantiate an event-conditioned asynchronous fusion setting in which lag is part of both sample construction and evaluation.}
  \dataset{} exposes $\tau_{\mathrm{lag}}$ explicitly across synthetic and real-news corpora, enabling lag-stratified analysis and freshness-aware evaluation under a causal walk-forward protocol.
  The resulting corpus reveals a non-monotonic utility pattern, with the strongest web signal appearing in the 30--60 minute post-publication range.
  \item \textbf{We propose a design principle for asynchronous fusion that separates grounding from trust control, and we instantiate it in \CGCMA{}.}
  Text first attends over the dense price stream to identify event-relevant market states, after which a learned per-dimension gate decides how much of that grounded context should be injected based on agreement, web features, and lag.
  \item \textbf{We provide cost-aware evidence for grounded conditional fusion on a real-news financial stress test.}
  Under a 5 bps cost-aware rolling protocol with strong price-only and multimodal baselines, \CGCMA{} achieves the strongest mean downstream Sharpe in a two-seed experiment batch ($+1.047\pm0.793$).
  Matched seed-fold comparisons show a stable margin over the price-only Transformer, while the margin over TimesNet-lite is positive but not statistically separated at the current sample scale.
  We therefore position the result as evidence for cost-aware asynchronous fusion rather than as a definitive market-wide trading claim.
\end{enumerate}

The remainder of this paper is organized as follows.
Section~\ref{sec:related} situates this paper in the literature.
Section~\ref{sec:problem} formalizes asynchronous multimodal market prediction.
Section~\ref{sec:method} describes the \CGCMA{} architecture and its asynchronous input structure.
Sections~\ref{sec:exp} and~\ref{sec:results} present the evaluation protocol and results.
Sections~\ref{sec:limitations}--\ref{sec:conclusion} discuss limitations, implications, and conclusions.
